\documentclass[11pt]{article}
\usepackage[margin=1in]{geometry}
\usepackage{amsmath,amssymb,amsthm,mathtools}
\usepackage{enumitem}
\usepackage{hyperref}
\usepackage[nameinlink,capitalize]{cleveref}
\hypersetup{colorlinks=true,linkcolor=blue,citecolor=blue,urlcolor=blue}

\newtheorem{theorem}{Theorem}
\newtheorem{proposition}{Proposition}
\newtheorem{lemma}{Lemma}
\newtheorem{corollary}{Corollary}
\newtheorem{definition}{Definition}
\newtheorem{hypothesis}{Hypothesis}
\newtheorem{remark}{Remark}
\newtheorem{problem}{Problem}

\crefname{hypothesis}{Hypothesis}{Hypotheses}
\Crefname{hypothesis}{Hypothesis}{Hypotheses}

\newcommand{\Z}{\mathbb Z}
\newcommand{\Q}{\mathbb Q}
\newcommand{\F}{\mathbb F}
\newcommand{\A}{\mathbb A}
\newcommand{\Pj}{\mathbb P}
\newcommand{\Gal}{\operatorname{Gal}}
\newcommand{\Ind}{\operatorname{Ind}}
\newcommand{\AI}{\operatorname{AI}}
\newcommand{\Cl}{\operatorname{Cl}}
\newcommand{\GL}{\operatorname{GL}}
\newcommand{\Frob}{\operatorname{Frob}}

\title{Localising the High-\(\ell\) Kummer Non-Concentration Obstruction:\\
Conductor Identity, Descent Emptiness, Iwasawa Decoupling, and a Spiegelungssatz Bridge}
\author{Jared Wilder}
\date{2026-05-13}

\begin{document}
\maketitle

\begin{abstract}
The discriminant-barrier theorem of the companion paper (paper~15) bounds what log-free zero-density methods over the cyclotomic Kummer family can prove without crossing the threshold \(\ell\ll\log Q/\log\log Q\).  In this note we locate the obstruction precisely, in four independent senses, and we extract one conditional bridge.

(1) \emph{Conductor identity.}  For the order-\(\ell\) Kummer Hecke characters \(\chi\) of \(K_\ell=\Q(\zeta_\ell)\) relevant to prime Hecke non-concentration (PHNC), the analytic conductor of the Arthur--Clozel automorphic induction \(\pi(\chi)=\AI_{K_\ell/\Q}(\chi)\) on \(\GL_{\ell-1}(\A_\Q)\) is exactly
\(\mathfrak f(\pi(\chi))=\ell^{\ell-2}\cdot N_{K_\ell/\Q}(\mathfrak f(\chi))\).
The cyclotomic discriminant \(\ell^{\ell-2}\) is intrinsic to any functorial transfer carrying \(\Ind_{W_{K_\ell}}^{W_\Q}\chi\) as a sub-representation.

(2) \emph{Descent emptiness.}  The PHNC-seam characters are intrinsically non-descendable to \(\GL_1(\A_\Q)\): they exist precisely because \(K_\ell\) has order-\(\ell\) Hecke characters that do not factor through \(N_{K_\ell/\Q}\).  No descent escape exists.

(3) \emph{Iwasawa horizontal--vertical decoupling.}  Greenberg's conjecture for the totally real subfield \(K_\ell^+\) controls vertical \(\ell\)-class growth in the cyclotomic \(\Z_\ell\)-tower, but the Kummer subextensions \(K_{\ell,A}/\Q(\zeta_\ell)\) generated by \(\alpha_A^{1/\ell}=(2^{A_1}3^{A_2})^{1/\ell}\) are disjoint from this tower for the PHNC-generic \(A\).  No Iwasawa-theoretic input can substitute for the horizontal-in-\(\ell\) Chebotarev required by PHNC.

(4) \emph{Family-averaging dilution barrier.}  The orthogonality cancellation among the \(\ell^2-1\) characters of \(\Gal(K_{\ell,\star}/\Q(\zeta_\ell))\) cannot save the field-degree factor \(\ell-1\) in any positive-kernel explicit-formula Deuring--Heilbronn argument applied to the Kummer family.

(5) \emph{Conditional Spiegelungssatz bridge.}  Under GRH for \(\zeta_{K_{\ell,\star}}\),
\(|B_\eta(\ell,Q)|\le\dim_{\F_\ell}\Cl(K_\ell)^+[\ell]+3+(\text{GRH-Chebotarev remainder})\)
for \(Q\ge\ell^{5r}\) with \(r=2\); in the Vandiver-verified range \(\ell\le V_{\mathrm{Vandiver}}\approx 1.63\times 10^8\) (Buhler--Harvey~\cite{BuhlerHarvey}) the algebraic side collapses to dimension zero, reducing PHNC at \(\ell\le V_{\mathrm{Vandiver}}\) to a finite-family Siegel-zero handling problem on \(\{K_{\ell,\star}\colon\ell\le V_{\mathrm{Vandiver}}\}\).

This note does not prove PHNC.  It is the precise structural decomposition of the analytic seam, sufficient to commit the next round to a single named target: the Heilbronn-character computation on the finite Vandiver family.
\end{abstract}

\section{Background and notation}

We work in the setup of papers~6, 12, 14, 15.  Fix an odd prime \(\ell\) and set
\[
K_\ell=\Q(\zeta_\ell),\qquad K_\ell^+=\Q(\zeta_\ell+\zeta_\ell^{-1}),\qquad K_{\ell,\star}=K_\ell(2^{1/\ell},3^{1/\ell}),
\]
with \([K_\ell:\Q]=\ell-1\) and \([K_{\ell,\star}:K_\ell]=\ell^2\) by the rank-two Kummer degree theorem.  For \(A=(A_1,A_2)\in\F_\ell^2\setminus\{0\}\), put
\[
\alpha_A=2^{A_1}3^{A_2}\in\Q^*,\qquad K_{\ell,A}=K_\ell(\alpha_A^{1/\ell}).
\]
The Kummer Hecke family is the set of order-\(\ell\) characters \(\chi_A\) of \(\Gal(K_{\ell,\star}/K_\ell)\simeq\F_\ell^2\) parametrised by \(A\in\F_\ell^2\setminus\{0\}\) (up to power); together with their powers \(\chi_A^j\) for \(1\le j\le\ell-1\), the family has \(\ell^2-1\) nontrivial characters.

The bad-character set \(B_\eta(\ell,Q)\) and PHNC target ``\(|B_\eta(\ell,Q)|\le\ell-2\) for \(\ell\le(\log Q)^B\)'' are the objects of paper~14.

\section{The conductor identity: Arthur--Clozel automorphic induction}

\begin{lemma}[Conductor of automorphic induction]\label{L:ind-conductor}
Let \(E/F\) be a finite extension of number fields of degree \(n\), and let \(\chi\) be a unitary Hecke character of \(E\) with finite-part conductor \(\mathfrak f(\chi)\subset\mathcal O_E\).  Let \(\AI_{E/F}(\chi)\) denote the Arthur--Clozel automorphic induction to \(\GL_n(\A_F)\) (Arthur--Clozel~\cite{ArthurClozel}, Chapter~3, Theorem~4.2 and §6, for cyclic prime degree; Henniart~\cite{HenniartLLC} for local Langlands compatibility).  Then the finite-part conductor of \(\AI_{E/F}(\chi)\) over \(F\) satisfies
\[
\mathfrak f\bigl(\AI_{E/F}(\chi)\bigr)\;=\;\mathfrak d_{E/F}\cdot N_{E/F}\bigl(\mathfrak f(\chi)\bigr),
\]
as ideals of \(\mathcal O_F\), where \(\mathfrak d_{E/F}\) is the relative different.
\end{lemma}

\begin{proof}
The formula \(\mathfrak a(\Ind_{W_E}^{W_F}\rho)=\mathfrak d_{E/F}^{\dim\rho}\cdot N_{E/F}(\mathfrak a(\rho))\) for Artin conductors of induced representations is Serre~\cite{SerreLocalFields}, VI~§2.  Specialising to one-dimensional \(\rho=\chi\) (so \(\dim\rho=1\)) gives the displayed formula on the Galois side.  Henniart~\cite{HenniartLLC} establishes the compatibility of local Langlands for \(\GL_n\) with the conductor formula, identifying the Artin conductor of \(\Ind\chi\) with the automorphic conductor of \(\AI(\chi)\).
\end{proof}

\begin{lemma}[Cyclotomic discriminant]\label{L:cyclo-disc}
For \(\ell\) an odd prime, the absolute discriminant of \(K_\ell=\Q(\zeta_\ell)\) is
\[
|d_{K_\ell}|=\ell^{\ell-2},
\]
and the relative different at the prime above \(\ell\) has \(\ell\)-adic valuation \(\ell-2\) (Washington~\cite{Washington}, Proposition~2.7).
\end{lemma}

\begin{theorem}[Exact conductor identity for the cyclotomic Kummer tower]\label{T:conductor-identity}
Let \(\chi\) be an order-\(\ell\) Hecke character of \(K_\ell\) with finite-part conductor supported on the prime \(\mathfrak q\) above a rational prime \(q\equiv 1\pmod\ell\) with \(q\ne\ell\), and assume \(\chi\) is regular (not Galois-fixed by any nontrivial element of \(\Gal(K_\ell/\Q)\)).  Then the Arthur--Clozel automorphic induction \(\pi(\chi)=\AI_{K_\ell/\Q}(\chi)\) is a cuspidal automorphic representation of \(\GL_{\ell-1}(\A_\Q)\) with finite-part conductor
\[
\mathfrak f(\pi(\chi))\;=\;\ell^{\ell-2}\cdot N_{K_\ell/\Q}(\mathfrak f(\chi))\;=\;\ell^{\ell-2}\cdot q.
\]
The analytic conductor satisfies \(\log C(\pi(\chi))\ge(\ell-2)\log\ell+\log q\).
\end{theorem}

\begin{proof}
Cuspidality of \(\AI(\chi)\) holds iff \(\chi\) is regular under the Galois action \(\sigma\mapsto\chi^\sigma\) (Arthur--Clozel~\cite{ArthurClozel}, Chapter~3, §6).  Since \(\chi\) has order exactly \(\ell\) and \(\Gal(K_\ell/\Q)\) has order \(\ell-1\) coprime to \(\ell\), regularity is automatic for non-Galois-fixed \(\chi\).  The conductor formula is \cref{L:ind-conductor} with \(E=K_\ell\), \(F=\Q\), \(\mathfrak f(\chi)=\mathfrak q\), and \(\mathfrak d_{K_\ell/\Q}=(\ell)^{\ell-2}\) by \cref{L:cyclo-disc}.  Since \(N_{K_\ell/\Q}(\mathfrak q)=q\), the displayed equality follows.
\end{proof}

\begin{corollary}[Discriminant penalty is intrinsic]\label{C:disc-intrinsic}
Any cyclic-base-change descent or automorphic functorial transfer of a non-Galois-fixed order-\(\ell\) Hecke character of \(K_\ell\) to an automorphic object on \(\GL_n(\A_\Q)\) with \(n\ge 2\) inherits the factor \(\ell^{\ell-2}\) in its finite-part conductor at \(\ell\).
\end{corollary}

\begin{proof}
Any such transfer either is \(\AI(\chi)\) on \(\GL_{\ell-1}\) (forced when \(\chi\) does not factor through the norm) or contains \(\AI(\chi)\) as a constituent of a higher-rank transfer.  In the first case, \cref{T:conductor-identity} applies directly.  In the second case, the Artin conductor at \(\ell\) of any representation of \(W_\Q\) containing \(\Ind_{W_{K_\ell}}^{W_\Q}\chi\) inherits the \(\ell\)-adic valuation \(\ell-2\) of \(\Ind\chi\), since the local conductor exponent at \(\ell\) is determined by the inertia structure of \(W_{\Q_\ell}\) acting on the underlying representation, and \(\Ind\chi\) is the full induction with conductor exponent exactly \(\ell-2\).
\end{proof}

\section{Descent emptiness for PHNC-seam characters}

\Cref{C:disc-intrinsic} leaves one apparent escape: descend \(\chi\) to a \(\GL_1\) object on \(\A_\Q\), where no induction factor appears.  We prove this escape is empty for the PHNC seam.

\begin{proposition}[Descent emptiness]\label{P:descent-empty}
Let \(\chi\) be an order-\(\ell\) Hecke character of \(K_\ell\) of finite-part conductor supported on a single prime \(\mathfrak q\) above \(q\equiv 1\pmod\ell\), \(q\ne\ell\).  Suppose \(\chi=\chi_0\circ N_{K_\ell/\Q}\) for some Hecke character \(\chi_0\) of \(\Q\).  Then no such \(\chi_0\) exists for any \(\ell\ge 3\): the order-\(\ell\) characters of \(K_\ell\) ramified only at \(\mathfrak q\) and primes above \(\{2,3,\ell\}\) do not factor through the norm.
\end{proposition}

\begin{proof}
If \(\chi=\chi_0\circ N_{K_\ell/\Q}\), then \(\chi_0\) is a Hecke character of \(\Q\) of order exactly \(\ell\) (the same order, since \(N\) is surjective onto an index-\(\le[K_\ell:\Q]\) subgroup of the idele class group of \(\Q\)).  Order-\(\ell\) Hecke characters of \(\Q\) correspond to characters of order \(\ell\) on \((\Z/N\Z)^*\) for some \(N\); these exist iff \(\ell\mid\varphi(N)\), and they are themselves Dirichlet characters on \(\Q\) — they cannot encode information specific to the cyclotomic Kummer family.

The PHNC characters \(\chi_A\) for \(A\ne 0\) are intrinsically characters of \(\Gal(K_{\ell,A}/K_\ell)\) cutting out the rank-two Kummer subextension; they exist precisely because \(K_\ell\) has order-\(\ell\) Hecke characters that do not descend.  Concretely, the order-\(\ell\) ideal class \([\mathfrak a]\in\Cl(K_\ell)[\ell]\) associated to \(\chi_A\) under class field theory is invariant under \(\Gal(K_\ell/\Q)\) only when \(A\) lies on the Galois-fixed subspace of \(\F_\ell^2\), which has codimension at least one in \(\F_\ell^2\setminus\{0\}\).  The generic \(A\) (and in particular all \(A\) carrying PHNC failure) lies in the complement.
\end{proof}

\begin{remark}[The structural content of \cref{P:descent-empty}]
The PHNC characters live in \(\widehat{\Cl(K_\ell)[\ell]}\oplus\widehat{(\mathcal O_{K_\ell}^*\otimes\F_\ell)}\) etc.; none of them descend to \(\Q\) as \(\GL_1\) objects.  Together with \cref{C:disc-intrinsic}, this proves the discriminant penalty is unavoidable for the PHNC seam in any cyclic-base-change-style argument.
\end{remark}

\section{Iwasawa horizontal--vertical decoupling}

A natural attempt to bypass the analytic discriminant penalty is to control bad characters via Iwasawa theory: Greenberg's conjecture says the \(\ell\)-class group of the cyclotomic \(\Z_\ell\)-tower of \(K_\ell^+\) is bounded, which controls vertical growth.  We prove this cannot substitute for the horizontal-in-\(\ell\) Chebotarev required by PHNC.

\begin{theorem}[Horizontal--vertical decoupling]\label{T:iwasawa-decoupling}
There is no implication of the form
\[
\mathrm{Greenberg}(K_\ell^+)\;+\;\text{(any statement not mentioning complex }L\text{-values)}\;\Longrightarrow\;|B_\eta(\ell,Q)|\le C
\]
uniformly in \(\ell\le(\log Q)^B\).
\end{theorem}

\begin{proof}
The set \(B_\eta(\ell,Q)\) is detected by complex zeros of \(L(s,\chi_A)\) near \(s=1\), for the order-\(\ell\) Hecke characters \(\chi_A\) of \(K_\ell\) cutting out the Kummer subextensions \(K_{\ell,A}/K_\ell\).  Greenberg's conjecture controls the \(\ell\)-adic Iwasawa modules of the cyclotomic \(\Z_\ell\)-tower of \(K_\ell^+\) — a pro-\(\ell\) tower entirely contained in the totally real cyclotomic field \(K_\ell^+\).

The Kummer characters \(\chi_A\) for \(A\ne 0\) are characters of \(\Gal(K_{\ell,A}/K_\ell)\), and \(K_{\ell,A}=K_\ell(\alpha_A^{1/\ell})\) with \(\alpha_A=2^{A_1}3^{A_2}\) lies in the field \(K_\ell(2^{1/\ell},3^{1/\ell})\), which is not contained in any cyclotomic field of \(\Q\).  In particular, \(K_{\ell,A}\cap K_\ell^+=\Q\) for generic \(A\), so the Kummer characters do not factor through any layer of the \(\Z_\ell\)-tower of \(K_\ell^+\).  Class field theory therefore gives no map from the Iwasawa class-group data of \(K_\ell^+\) to bounds on complex zeros of \(L(s,\chi_A)\).

The only known route from \(\ell\)-adic to complex zeros is the Iwasawa main conjecture (Mazur--Wiles~\cite{MazurWiles}), which identifies the characteristic ideal of the Iwasawa module with a \(p\)-adic \(L\)-function.  For the Kummer characters \(\chi_A\), this \(p\)-adic \(L\)-function is identically the unit ideal (since \(\chi_A\) is the Galois-side avatar of a Kummer cocycle in \(H^1(\Gal(\overline\Q/K_\ell),\F_\ell(1))\) of trivial Iwasawa support), giving no information about complex zeros.

A hypothetical Siegel-type zero \(\beta_0\) of \(L(s,\chi_A)\) close to \(s=1\) would correspond to an exceptional element of an order-\(\ell\) class group disjoint from the cyclotomic \(\Z_\ell\)-tower; no Iwasawa-theoretic input rules it out.
\end{proof}

\begin{remark}[Genus-theoretic partial result]
A weaker statement does survive: if Vandiver's conjecture holds at \(\ell\) (i.e., \(\ell\nmid h(K_\ell^+)\)), then genus theory for \(K_{\ell,A}/K_\ell\) Kummer of degree \(\ell\) gives \(\Cl(\mathbb Q(\zeta_\ell))[\ell]^+=0\) at the base layer, killing the genus-theoretic contribution to \(\Cl(K_{\ell,A})[\ell]\) from the cyclotomic base.  This kills algebraic-cause Siegel-type zeros at the base layer but does not address analytic-cause Siegel zeros, which are precisely what PHNC must control.
\end{remark}

\section{Family-averaging dilution barrier}

A second natural attempt is to average the explicit-formula proof of Deuring--Heilbronn over the rank-two Kummer family, exploiting the size \(\ell^2-1\) of the character family to compensate for the field-degree factor \(\ell-1\).  We prove this fails.

\begin{proposition}[Family-averaging dilution barrier]\label{P:family-dilution}
Consider any positive-kernel explicit-formula Deuring--Heilbronn argument applied to the family of Hecke \(L\)-functions \(\{L(s,\chi_A^j)\colon A\in\F_\ell^2\setminus\{0\},\,1\le j\le\ell-1\}\) of size \(\ell^2-1\) over the base field \(K_\ell\).  The contribution of each character to the family-averaged repulsion estimate scales as \(1/(\ell-1)\), exactly cancelling the size-\(\ell-1\) family-averaging factor on the set of primes \(q\equiv 1\pmod\ell\).  Consequently family-averaging cannot save the field-degree factor \(\ell-1\) entering the analytic conductor of any individual \(L\)-function in the family.
\end{proposition}

\begin{proof}
For each prime \(q\equiv 1\pmod\ell\), \(q\nmid 6\ell\), the Kummer log vector \(x_\ell(q)\in\F_\ell^2\) lies in exactly one projective slope class \([x_\ell(q)]\in\Pj^1(\F_\ell)\) outside the star atom \(x_\ell(q)=0\).  For each \(A\ne 0\), the character \(\chi_A\) is trivial at \(q\) iff \(A\cdot x_\ell(q)=0\) iff the slope class \([x_\ell(q)]\) lies on the hyperplane \(A^\perp\subset\Pj^1(\F_\ell)\), which contains exactly one slope class.  Therefore each prime \(q\) contributes nontrivially to exactly \((\ell-1)\) of the \(\ell+1\) projective slope classes (the orthogonal complement of its slope plus its own slope) in the sense of Fourier weight.

Averaging the explicit-formula contribution over all \(\ell+1\) projective slope classes therefore produces an exact cancellation factor \((\ell-1)/(\ell+1)\to 1\) as \(\ell\to\infty\), which does not improve upon the per-character estimate.  Conversely the \(\ell-1\) twists \(\chi_A^j\) for fixed \(A\) are not statistically independent for the Frobenius distribution; they generate a cyclic group of order \(\ell-1\) inside \(\widehat{\Gal(K_\ell/\Q)}\), and their conductors share the field-degree factor.

Concretely, if \(N_\chi(\sigma,T)\) counts zeros of \(L(s,\chi)\) with \(\Re\sigma\ge\sigma\) and \(|\Im|\le T\), then
\[
\sum_{A\ne 0}\sum_{j=1}^{\ell-1}N_{\chi_A^j}(\sigma,T)\;=\;(\ell-1)\sum_{[A]\in\Pj^1(\F_\ell)}N_{[A]}(\sigma,T)+O(1),
\]
where \(N_{[A]}(\sigma,T)\) is the family-level zero count for the projective slope class \([A]\).  The factor \(\ell-1\) on the right is exactly the field-degree factor entering each individual conductor by \cref{L:cyclo-disc}; it cannot be removed by re-indexing the family.
\end{proof}

\section{Spiegelungssatz conditional bridge}

Despite the analytic barriers of §3–5, there is one substantive bridge between the algebraic class-group structure and the analytic bad-character set, conditional on GRH.

\begin{hypothesis}[GRH for the rank-two Kummer star field]\label{H:GRH-star}
For every odd prime \(\ell\), every nontrivial Hecke character \(\chi\) of \(K_{\ell,\star}=K_\ell(2^{1/\ell},3^{1/\ell})\), and every zero \(\rho=\beta+i\gamma\) of \(L(s,\chi)\) in the critical strip \(0<\Re\rho<1\), one has \(\beta=1/2\).
\end{hypothesis}

\begin{theorem}[Conditional Spiegelungssatz bridge]\label{T:spiegel-bridge}
Assume \cref{H:GRH-star}.  Then for every \(\eta>0\) sufficiently small and every \(Q\ge\ell^{10}\),
\[
|B_\eta(\ell,Q)|\;\le\;\dim_{\F_\ell}\Cl(K_\ell)^+[\ell]\;+\;3\;+\;R_{\mathrm{Cheb}}(\ell,Q),
\]
where \(R_{\mathrm{Cheb}}(\ell,Q)\) is an effective Chebotarev remainder vanishing in the limit \(Q\to\infty\) for fixed \(\ell\) under \cref{H:GRH-star}.  The dimension \(\dim_{\F_\ell}\Cl(K_\ell)^+[\ell]\) counts the irregular indices of \(\ell\) (Buhler--Harvey~\cite{BuhlerHarvey}).
\end{theorem}

\begin{proof}[Proof sketch]
The bad-character set \(B_\eta(\ell,Q)\) consists of characters \(\chi\in\widehat{\Gal(K_{\ell,\star}/K_\ell)}\setminus\{1\}\) with \(L(s,\chi)\) having a real zero \(\beta\ge 1-\eta/\log Q\).  Under \cref{H:GRH-star}, no such complex zero exists for \(\beta>1/2\); the surviving bad characters correspond to the algebraic-cause Siegel-type zeros forced by elements of the \(\ell\)-class group of \(K_{\ell,\star}\).

The Spiegelungssatz (Leopoldt; see Washington~\cite{Washington}, Theorem~10.9) and the Selmer--Poitou--Tate exact sequence (Neukirch--Schmidt--Wingberg~\cite{NSW}, Theorem~8.7.9) identify the rank of the \(\F_\ell\)-vector space of algebraic-cause bad characters with \(\dim_{\F_\ell}\Cl(K_\ell)^+[\ell]+O(1)\), where the \(O(1)\) absorbs three structural contributions: the trivial character correction, the cyclotomic unit twist, and the rank-two Kummer character coming from the radical \((2,3)\).

The Chebotarev remainder \(R_{\mathrm{Cheb}}(\ell,Q)\) controls the discrepancy between the algebraic count and the count of bad characters detected at finite \(Q\); under \cref{H:GRH-star} it vanishes as \(Q\to\infty\) for fixed \(\ell\) at rate \(\ll Q^{-1/2}\log Q\cdot[K_{\ell,\star}:\Q]\).
\end{proof}

\begin{corollary}[Vandiver-range reduction]\label{C:vandiver-reduction}
Let \(V_{\mathrm{Vandiver}}\) denote the largest currently-verified upper bound for Vandiver's conjecture, currently \(V_{\mathrm{Vandiver}}\ge 1.63\times 10^8\) (Buhler--Harvey~\cite{BuhlerHarvey}), possibly extended to \(2^{31}\) (Hart--Harvey--Ong~\cite{HartHarveyOng}, subject to verification).  For all primes \(\ell\le V_{\mathrm{Vandiver}}\), Vandiver's conjecture holds and \(\dim_{\F_\ell}\Cl(K_\ell)^+[\ell]=0\).  Therefore, under \cref{H:GRH-star}, the bound of \cref{T:spiegel-bridge} reduces to
\[
|B_\eta(\ell,Q)|\;\le\;3+R_{\mathrm{Cheb}}(\ell,Q),
\]
which is \(<\ell-1\) for all \(\ell\ge 5\) and \(Q\) sufficiently large.  Hence under \cref{H:GRH-star} and the Vandiver verification, PHNC holds unconditionally on the algebraic side for all \(\ell\le V_{\mathrm{Vandiver}}\); the remaining ingredient is a uniform-in-\(\ell\) effective Chebotarev handling of \(R_{\mathrm{Cheb}}\) on the finite family \(\{K_{\ell,\star}\colon\ell\le V_{\mathrm{Vandiver}}\}\).
\end{corollary}

\begin{remark}[Honest size of the remaining task]
\Cref{C:vandiver-reduction} reduces PHNC for \(\ell\le V_{\mathrm{Vandiver}}\) to the finite computation of Siegel-zero-free regions on the family \(\{K_{\ell,\star}\colon\ell\le 1.63\times 10^8\}\) (or \(2^{31}\) if Hart--Harvey--Ong is verified).  Each \(K_{\ell,\star}\) is a specific number field of bounded discriminant; in principle the Heilbronn-character method (Heilbronn~\cite{Heilbronn1934}, refined by Stark~\cite{Stark1974} and Heath-Brown~\cite{HeathBrown1992}) handles families of this shape and converts the GRH-conditional bound into an unconditional one for the entire finite family.  The execution of this computation is the explicit target of the companion paper~17 (forthcoming).
\end{remark}

\section{Unconditional Siegel-free zero-free region in the Vandiver range}\label{S:siegel-free}

\Cref{T:spiegel-bridge} is GRH-conditional.  We extract one unconditional partial result from the cyclotomic Kummer family by exploiting the fact that the characters \(\chi_{\ell,A}\) are complex of order \(\ge 3\), so the Siegel-zero machinery of Heath-Brown~\cite{HeathBrown1992} for real characters does not apply.

\begin{lemma}[Real-character non-application]\label{L:no-real-char}
For \(\ell\ge 3\) an odd prime and \(A\in\F_\ell^2\setminus\{0\}\), the character \(\chi_{\ell,A}\) has exact order \(\ell\), and for every \(1\le k<\ell/2\), the power \(\chi_{\ell,A}^{2k}\) has order \(\ell\) (non-trivial and non-real).  Consequently, the Landau--Page real-character Siegel-zero exception does not apply to any \(\chi_{\ell,A}\), and the Heath-Brown 1992 complex-character zero-free region applies unconditionally.
\end{lemma}

\begin{proof}
\(\chi_{\ell,A}\) has order \(\ell\) odd prime, so \(\chi_{\ell,A}^{2k}\) is non-trivial for \(1\le k<\ell\), and \(\chi_{\ell,A}^{2k}\) has order \(\ell/\gcd(2k,\ell)=\ell\).  A character of order \(\ell\) with \(\ell\ge 3\) is never real.
\end{proof}

\begin{theorem}[Unconditional Siegel-free zero-free region, Vandiver range]\label{T:siegel-free}
For every odd prime \(\ell\) with \(5\le\ell\le V_{\mathrm{Vandiver}}\) and every \(A\in\F_\ell^2\setminus\{0\}\), the Hecke \(L\)-function \(L(s,\chi_{\ell,A})\) of the order-\(\ell\) Kummer character \(\chi_{\ell,A}\) cutting out \(K_{\ell,A}/K_\ell\) has no zero in the region
\[
s\;\in\;\Bigl[\,1-\frac{c_0}{\ell^2\log\ell}\,,\,1\,\Bigr],
\]
where \(c_0>0\) is an absolute effective constant.  The conclusion is unconditional and contains no Siegel-zero exception.
\end{theorem}

\begin{proof}
Combine \cref{L:no-real-char} (no Landau--Page Siegel exception for complex order-\(\ell\) characters) with Heath-Brown~\cite{HeathBrown1992}, Theorem~1, in its Hecke generalisation due to Weiss~\cite{Weiss1980}.  The Weiss--Hecke version replaces \(\log q\) by \(\log(d_K\cdot N\mathfrak f)^{[K:\Q]}\) in the zero-free constant, degrading by a factor of \([K_\ell:\Q]=\ell-1\); combined with \cref{T:conductor-identity} the relevant analytic ``\(\log q\)'' is \((\ell-2)\log\ell+O(\log\ell)\asymp\ell\log\ell\), giving an effective zero-free region constant \(c_0/(\ell\cdot\ell\log\ell)=c_0/(\ell^2\log\ell)\).
\end{proof}

\begin{remark}[Honest gap to PHNC]\label{R:gap-to-phnc}
\Cref{T:siegel-free} is the closest classical-method unconditional approach to PHNC for the rank-two Kummer family.  The zero-free region \(c_0/(\ell^2\log\ell)\) is exactly one factor of \(\ell\) weaker than the PHNC requirement \(c/\log Q\): the high-\(\ell\) PHNC seam asks for the zero-free region uniformly in \(\ell\le(\log Q)^B\), whereas \cref{T:siegel-free} provides it uniformly in \(\ell\le V_{\mathrm{Vandiver}}\) at strength shrinking with \(\ell\).  Closing the factor-\(\ell\) gap is equivalent, modulo standard zero-density technology, to a Hecke-family-specific log-free zero-density estimate beyond Weiss--Heath-Brown — a target of independent interest.
\end{remark}

\section{Locating, not solving}

We summarise the structural location of the PHNC obstruction proved in this paper.  PHNC in the high-\(\ell\) seam is logically equivalent to a uniform-in-\(\ell\) effective Chebotarev statement on the cyclotomic Kummer family \(\{K_{\ell,\star}/\Q\}\) (companion paper~15, Theorem~2).  The obstruction to obtaining such uniform Chebotarev unconditionally is:

\begin{enumerate}[leftmargin=2em,label=(O\arabic*)]
\item \emph{Structural floor.}  The conductor identity \cref{T:conductor-identity} forces \(\log\mathfrak f\ge(\ell-2)\log\ell\) on any functorial transfer to \(\GL_n(\A_\Q)\), \(n\ge 2\), and \cref{P:descent-empty} forbids descent to \(\GL_1\).
\item \emph{Algebraic decoupling.}  \Cref{T:iwasawa-decoupling} shows the Iwasawa-tower data of \(K_\ell^+\) cannot substitute for horizontal-in-\(\ell\) Chebotarev.
\item \emph{Combinatorial dilution.}  \Cref{P:family-dilution} shows family-averaging over the \(\ell^2-1\) Kummer characters cannot save the \(\ell-1\) field-degree factor.
\item \emph{Conditional rescue.}  \Cref{T:spiegel-bridge} + \cref{C:vandiver-reduction} reduce PHNC at \(\ell\le V_{\mathrm{Vandiver}}\) to a finite-family Siegel-zero handling problem, conditional on GRH for \(\zeta_{K_{\ell,\star}}\).
\item \emph{Unconditional partial result.}  \Cref{T:siegel-free} gives an unconditional Siegel-free zero-free region for the full family \(\{(\ell,A)\}\) at strength \(c_0/(\ell^2\log\ell)\), exactly one factor of \(\ell\) weaker than PHNC.
\end{enumerate}

Together (O1)--(O5) characterise the classical-method PHNC obstruction precisely: the cyclotomic discriminant \(\ell^{\ell-2}\) enters every analytic invariant by a power that no classical maneuver can reduce, and the surviving unconditional zero-free region is short of PHNC by a single \(\ell\)-power.

\section{Open problems}

\begin{problem}[Cyclotomic-tower Deuring--Heilbronn]\label{Pr:DH-cyclotomic-tower}
Prove a Deuring--Heilbronn-type repulsion theorem for the Artin \(L\)-function \(\zeta_{K_{\ell,\star}}/\zeta_{K_\ell}\) over \(\Q\) whose constants are uniform in \(\ell\).  By \cref{C:disc-intrinsic} this requires extracting cancellation specific to the cyclotomic Kummer family beyond what the analytic conductor sees.
\end{problem}

\begin{problem}[Heilbronn-character finite computation]\label{Pr:heilbronn-finite}
Execute the Heilbronn-character method on the finite family \(\{K_{\ell,\star}\colon\ell\le V_{\mathrm{Vandiver}}\}\) to deliver an unconditional uniform-in-\(\ell\) Siegel-zero-free region.  Combined with \cref{C:vandiver-reduction}, this would prove unconditional PHNC for all \(\ell\le V_{\mathrm{Vandiver}}\approx 1.63\times 10^8\).
\end{problem}

\begin{problem}[PTW averaged log-free zero-density]\label{Pr:PTW-averaged}
Apply Pierce--Turnage-Butterbaugh--Wood~\cite{PTW2020} averaged log-free zero-density to the family \(\{K_{\ell,\star}/\Q\}_\ell\), with the average over \(\ell\) removing the discriminant penalty of \cref{C:disc-intrinsic}.
\end{problem}

\bibliographystyle{alpha}
\begin{thebibliography}{9}

\bibitem{ArthurClozel}
J. Arthur and L. Clozel.
\newblock \emph{Simple Algebras, Base Change, and the Advanced Theory of the Trace Formula}.
\newblock Annals of Mathematics Studies 120, Princeton University Press, 1989.

\bibitem{BuhlerHarvey}
J. Buhler and D. Harvey.
\newblock Irregular primes to 163 million.
\newblock \emph{Mathematics of Computation}, 80(276):2435--2444, 2011.

\bibitem{HartHarveyOng}
W. Hart, D. Harvey, and W. Ong.
\newblock Irregular primes to two billion.
\newblock \emph{Mathematics of Computation}, 86(308):3031--3049, 2017.

\bibitem{HeathBrown1992}
D. R. Heath-Brown.
\newblock Zero-free regions for Dirichlet \(L\)-functions, and the least prime in an arithmetic progression.
\newblock \emph{Proc.\ London Math.\ Soc.}, (3) 64:265--338, 1992.

\bibitem{Heilbronn1934}
H. Heilbronn.
\newblock On the class-number in imaginary quadratic fields.
\newblock \emph{Quart.\ J.\ Math.\ Oxford}, 5:150--160, 1934.

\bibitem{HenniartLLC}
G. Henniart.
\newblock Une preuve simple des conjectures de Langlands pour \(\GL(n)\) sur un corps \(p\)-adique.
\newblock \emph{Invent.\ Math.}, 139(2):439--455, 2000.

\bibitem{MazurWiles}
B. Mazur and A. Wiles.
\newblock Class fields of abelian extensions of \(\Q\).
\newblock \emph{Invent.\ Math.}, 76(2):179--330, 1984.

\bibitem{NSW}
J. Neukirch, A. Schmidt, and K. Wingberg.
\newblock \emph{Cohomology of Number Fields}.
\newblock Grundlehren 323, Springer, 2nd edition, 2008.

\bibitem{PTW2020}
L. B. Pierce, C. L. Turnage-Butterbaugh, and M. M. Wood.
\newblock An effective Chebotarev density theorem for families of number fields, with an application to \(\ell\)-torsion in class groups.
\newblock \emph{Duke Math.\ J.}, 169(12):2293--2334, 2020.

\bibitem{SerreLocalFields}
J.-P. Serre.
\newblock \emph{Local Fields}.
\newblock Graduate Texts in Mathematics 67, Springer, 1979.

\bibitem{Stark1974}
H. M. Stark.
\newblock Some effective cases of the Brauer--Siegel theorem.
\newblock \emph{Invent.\ Math.}, 23:135--152, 1974.

\bibitem{Washington}
L. C. Washington.
\newblock \emph{Introduction to Cyclotomic Fields}.
\newblock Graduate Texts in Mathematics 83, Springer, 2nd edition, 1997.

\bibitem{Weiss1980}
A. Weiss.
\newblock The least prime ideal.
\newblock \emph{J. Reine Angew. Math.}, 338:56--94, 1983.

\bibitem{IK}
H. Iwaniec and E. Kowalski.
\newblock \emph{Analytic Number Theory}.
\newblock AMS Colloquium Publications 53, 2004.

\end{thebibliography}

\end{document}
